\chapter{Accumulation and the periodic law}

\textbf{The figures here are measured on the ideal filling and are deterministic.} The ideal filling
is defined first, because every later claim is a claim about it and not about a measured atom.

\section{The ideal filling}

The ledger at
\texttt{src/\allowbreak{}engine/\allowbreak{}python/\allowbreak{}representation/\allowbreak{}atom/\allowbreak{}element.py}
builds an element's electrons in Madelung order, lowest principal plus azimuthal first, and Hund's
rule inside a subshell, every magnetic orbital taking a spin-up electron before any takes spin-down.
The order fills 1s through 7p, and its capacities run to 118 exactly. It names every element from
hydrogen to oganesson and no further.

This is the necessary condition, the filling the periodic law would follow if nothing competed with
it. Real atoms deviate: chromium, copper and others fill against the order, and those measured
configurations belong to the oracle stage of the next chapter. The represent field carries the ideal,
and it says so. That no reader takes the ideal for a measurement.

\section{Exclusion is an accumulation}

Two fermions may not occupy one state. An electron added to an atom therefore cannot join the
electrons already in the lowest state; it takes the next state the exclusion leaves open, and the one
after that takes the next again. Electrons accumulate into successive states, and the stack of
occupied states is the shell structure. Exclusion is not a bookkeeping rule laid over the atom. It is
the reason the atom has structure at all.

That accumulation is faithful here for a reason the first chapter gave. A reading that rounded two
close states together would lose an electron from the stack and shorten a shell. The exact reading
never does, because its discrimination is unbounded, and an accumulation of fermions stays an
accumulation and never collapses. Exactness lets the engine hold the stack, and exclusion makes the
stack tall.

\section{The recurrence, against a drawn null}

The differentiating electron of an element is the last one the order adds, and its group signature is
its azimuthal type and its subshell's population with the principal number dropped. Dropping the
principal number leaves sodium and lithium reading the same signature, and that recurrence down the
table is the periodic law. It is measured by an equality match rate, how often a signature equals the
one a fixed number of elements ahead, in
\texttt{examples/\allowbreak{}particle\_physics/\allowbreak{}3\_reference/\allowbreak{}what\_exclusion\_builds.py}.

A match rate means nothing without the rate a shuffle of the same signatures reaches. Two
backgrounds are drawn, each deleting a property from the data. The first keeps the census and deletes
the order: a uniform permutation of the same signatures, which fixes every count exactly. Its floor
is the real floor, and the match rate at the row lengths falls to it. The recurrence is in the
order the shells fill and not in the counts. The second deletes the exclusion: every electron drops
to 1s, an element becomes a count of electrons in one state, every signature is distinct, and the
match rate and its floor are both zero. Without exclusion there is a ladder of counts and no table.
The real accumulation is the positive control, and it is the only arm that carries the recurrence.

\section{The row lengths, off the difference set}

The recurrence does not sit at one lag. The match rate stands above its floor at 32, at 18 and at 8,
and at their sums, because the rows of the table are 2, 8, 8, 18, 18, 32, 32 and not one length. A
reader reporting a single period would be wrong for most of the table, the harmonic trap the
periodicity primitive records from the other side. What is one number per row is the gap between
shell closures, and
\texttt{examples/\allowbreak{}particle\_physics/\allowbreak{}4\_measure/\allowbreak{}shell\_closures\_from\_the\_difference\_set.py}
reads them off the signatures: a closure is a filled p subshell, with helium closing the first row.

\begin{center}
\begin{tabular}{@{}lrrrrrrr@{}}
\toprule
closure at $Z$ & 2 & 10 & 18 & 36 & 54 & 86 & 118 \\
\midrule
row length     &   & 8  & 8  & 18 & 18 & 32 & 32 \\
\bottomrule
\end{tabular}
\end{center}

\noindent The closures are the noble gases, and the gaps between them are the row lengths, with no row
told to the reader. This is the difference-set reading of the crystallography chapter moved to one
dimension: a period is a difference between two things that agree. The differences are the complete
candidate set.
